\subsection{Adam 的更新规则}

Adam 的核心是同时维护梯度的一阶矩估计（均值）和二阶矩估计（未中心化方差），并用偏差校正来补偿初始化时的偏差：

$$
m_t = \beta_1 m_{t-1} + (1 - \beta_1) g_t \qquad \text{（一阶矩/动量）}
$$
$$
v_t = \beta_2 v_{t-1} + (1 - \beta_2) g_t^2 \qquad \text{（二阶矩/自适应学习率）}
$$
$$
\hat{m}_t = \frac{m_t}{1 - \beta_1^t}, \quad \hat{v}_t = \frac{v_t}{1 - \beta_2^t} \qquad \text{（偏差校正）}
$$
$$
\theta_{t+1} = \theta_t - \eta \cdot \frac{\hat{m}_t}{\sqrt{\hat{v}_t} + \epsilon}
$$

其中 $\beta_1 = 0.9$, $\beta_2 = 0.999$, $\epsilon = 10^{-8}$ 是常见默认值。

\begin{knowledgebox}{AdamW vs Adam + L2 正则化——一个微妙但重要的区别}
经典 Adam + L2 正则化把权重衰减加在梯度上：
$$
g_t' = g_t + \lambda \theta_t, \qquad \theta_{t+1} = \theta_t - \eta \cdot \frac{\hat{m}_t(g_t')}{\sqrt{\hat{v}_t(g_t')} + \epsilon}
$$
问题是：权重衰减项 $\lambda \theta_t$ 也被 Adam 的自适应学习率缩放了。对于梯度较大的参数，衰减效果被削弱；对于梯度较小的参数，衰减效果被放大。这并不是我们想要的行为。

AdamW（Loshchilov \& Hutter, 2019）把权重衰减解耦出来：
$$
\theta_{t+1} = (1 - \eta \lambda) \theta_t - \eta \cdot \frac{\hat{m}_t}{\sqrt{\hat{v}_t} + \epsilon}
$$
这样权重衰减不受自适应学习率的影响，对所有参数施加均匀的正则化。

\textbf{实践结论}：现代 LLM 训练几乎都使用 AdamW 而非 Adam + L2。PyTorch 中 \texttt{torch.optim.AdamW} 实现的就是解耦版本。
\end{knowledgebox}

